% !TeX root = ../main.tex
%------------------------------------------------------------------------------
%  Groundwater and Hydrogeology
%  Sources: Goodman et al. (1965); Heuer (1995); El Tani (2003);
%  ITA guidance on drainage and waterproofing.
%------------------------------------------------------------------------------
\chapter{Groundwater and Hydrogeology}
\label{ch:hydrogeology}

Groundwater affects underground construction in four ways that are frequently
treated as one. It reduces effective stress and therefore strength; it flows
into the excavation and must be handled and discharged; it destabilises the
face, particularly in granular ground; and it loads the completed lining, in
some cases for the whole service life. These require different analyses and
different control measures, and a project that treats groundwater as a single
topic tends to address whichever aspect is most familiar. This chapter
separates them (the third is taken up with face support in
Section~\ref{sec:face-support}, where the water pressure at the face enters the
required support pressure) and then examines the decision that follows from them all:
whether the finished structure drains or resists the full head. That decision
is taken early, on hydrogeological grounds, and it constrains all later
design.

\section{Groundwater regime and flow}
\label{sec:groundwater-regime}

\subsection{Conceptual hydrogeological models}
\label{sec:hydrogeological-model}

A conceptual model states where water enters the system, where it leaves, what
transmits it and what impedes it. Aquifers and aquitards are defined by
contrast rather than in absolute terms, and the contrast that matters is
usually several orders of magnitude in permeability. Perched systems, confined
aquifers whose piezometric head stands well above the top of the aquifer
itself, and recharge boundaries such
as rivers all change the answer to the same excavation problem.

The distinction that governs analysis is between porous-medium flow and
fracture-controlled flow. In a granular soil or a well-connected porous rock,
flow occupies the whole cross-section and an equivalent continuum treatment is
sound. In a low-porosity fractured rock, flow is confined to a small number of
connected features, and the equivalent continuum predicts the average
inflow but not its distribution. Inflow into a
tunnel in such ground arrives at discrete points rather than along the whole
drive, at rates far above the average.

Recognising which regime applies is a geological judgement made from the rock
type, the structure and the fracture connectivity, and it should be made
explicitly before any inflow calculation is attempted.

\subsection{Permeability and effective stress}
\label{sec:permeability}

Permeability is measured over very different support volumes. A laboratory
permeameter tests a specimen deliberately selected to be intact, and therefore
measures matrix permeability, frequently irrelevant in fractured rock. A
packer or Lugeon test in a borehole measures a few metres of rock around the
hole, including whatever fractures it happens to intersect. The Lugeon unit
converts to a hydraulic conductivity through the factor given in
Appendix~\ref{app:constants}. A pumping test
integrates over tens or hundreds of metres and is the only routine measurement
whose support volume approaches the scale of the problem.

The consequence is that measured permeability tends to rise with the scale of
the test; in crystalline rock the rise from laboratory to borehole to regional
scale can amount to orders of magnitude~\citep{clauser1992permeability}. A
design based on single-borehole values can therefore underestimate inflow,
most seriously where the flow is carried by sparse features that few boreholes
intersect. Where inflow governs, a pumping test is
worth more than a large number of packer tests.

Effective stress links water back to the mechanics. Excavation changes pore
pressure as well as total stress, and the response depends on whether the
ground drains during the time available. Low-permeability clays respond
essentially undrained during construction and drain over subsequent years,
so short-term face stability and long-term lining load are different
calculations on the same section. Permeable ground drains as fast as it is
loaded and only the drained case exists. The intermediate case (partially
drained, with the degree of drainage depending on advance rate) is the most
awkward and the most common in weathered rock.

\subsection{Inflow prediction}
\label{sec:inflow-prediction}

Steady-state analytical solutions give inflow per unit length of tunnel as a
function of permeability, head and tunnel radius, with the classical
forms~\citep{goodman1965inflow,eltani2003circular} differing in their
treatment of the free surface and the boundary conditions. The simplest, for
a tunnel of radius $r_0$ whose axis lies at a depth $H$ below a water table
held at constant level, with $H$ large in comparison with $r_0$, is
\[
q = \frac{2\pi k_w H}{\ln\left(2H/r_0\right)} ,
\]
where $q$ is the inflow per unit length of tunnel and $k_w$ the hydraulic
conductivity of the ground.
Such solutions are simple, transparent and dimensionally correct, but they are only as
good as the permeability supplied to them (Figure~\ref{fig:inflow-goodman}).
They also give the eventual
steady rate. The inflow on first intersecting a water-bearing feature is
higher while storage drains, and pumping capacity has to be sized against
that transient peak as well.

\begin{figure}[htbp]
\centering
\datafigure[\textwidth]{ch07/inflow-goodman}
\caption{Steady inflow per metre of tunnel from $q = 2\pi k_w H/\ln(2H/r_0)$ \citep{goodman1965inflow}; left $r_0 = 3$~m, right $k_w = 10^{-7}$~m/s, with $H$ the depth of the axis below the water table. Inflow follows $k_w$ over orders of magnitude but changes by about a third for a threefold change in radius. Top axis: $1~\mathrm{Lu} \approx 1.3\times10^{-7}$~m/s.}
\label{fig:inflow-goodman}
\end{figure}

Two systematic corrections are needed in fractured rock. The first is for
scale, as above. The second follows from the discrete nature of the flow: a
solution assuming uniform permeability distributes the predicted inflow evenly
along the drive, whereas the actual inflow concentrates at fractures and faults.
The total may be approximately right while the peak, which determines pumping
capacity and whether the face can be worked, is badly wrong
(Figure~\ref{fig:inflow-discrete}).

\begin{figure}[htbp]
\centering
\datafigure[\textwidth]{ch07/inflow-discrete}
\caption{Inflow along a 500~m drive: equivalent continuum against 14 randomly placed features carrying 90 per cent of the same total (synthetic; $k_w = 10^{-7}$~m/s, $H = 50$~m, $r_0 = 3$~m). The totals agree at the end of the drive, but the largest feature delivers 80~L/min within a metre, some 150 times the uniform rate.}
\label{fig:inflow-discrete}
\end{figure}

Expectations of accuracy should be modest~\citep{heuer1995estimating}. Inflow
predictions that fall within a factor of two or three of the outturn are
good; order-of-magnitude errors are
common where the flow regime was misdiagnosed. Design should therefore be
based on a range with a stated basis, and pumping and treatment capacity sized
against the upper bound rather than the central estimate. Probabilistic
treatment along the alignment (combining the probability of intersecting a
water-bearing feature with the inflow expected if one is met) is more
informative than a single deterministic value.

\section{Design consequences of groundwater}
\label{sec:groundwater-consequences}

\subsection{Drained and tanked linings}
\label{sec:drained-vs-tanked}

The choice between a drained and a tanked lining is the governing groundwater
decision (Figure~\ref{fig:drained-tanked-lining}). A drained lining admits water to a
drainage layer and conveys it to a sump, so the lining carries little or no
water pressure and can be correspondingly thinner. That embodied-carbon saving
has to be set against the operational carbon of pumping over the design
life (Chapter~\ref{ch:sustainability}). In
exchange the owner accepts permanent drawdown, perpetual pumping, and a
drainage system that will scale with calcite and require access for flushing
over the whole design life. A tanked lining resists the full hydrostatic head,
which means a thicker and more heavily reinforced structure, an uplift check
for shallow sections, and a watertightness requirement that the construction
must achieve. The groundwater regime is left largely
undisturbed, however, and the operator has nothing to pump.

\begin{figure}[htbp]
\centering
\datafigure[\textwidth]{ch07/drained-tanked-lining}
\caption{Drained and tanked linings (schematic). Left: a drainage layer (dashed ring) and sidewall drains relieve the water pressure on the lining and draw the water table down over the tunnel (dotted: original level). Right: a membrane keeps the water out, and the lining carries the full hydrostatic head, rising from crown to invert.}
\label{fig:drained-tanked-lining}
\end{figure}

Neither is universally correct, and the choice cannot be made on structural
grounds alone. It is decided by the environmental acceptability of drawdown,
the head, the permeability, the availability of maintenance access, and the
owner's appetite for an operating liability in exchange for a capital saving.
Because it determines lining thickness, reinforcement, waterproofing system
and drainage detailing, it must be made early; reversing it at detailed design
is effectively a redesign.

Design head requires its own judgement. Water tables recover after
construction, sometimes decades later and sometimes above the pre-construction
level where regional abstraction has ceased. A tanked structure designed to
the head measured during construction may be under-designed for the head it
eventually sees, and shallow structures require an uplift check against the
recovered water table with the structure empty.

\subsection{Drawdown and environmental effects}
\label{sec:drawdown}

Drawdown causes settlement by increasing effective stress in compressible
strata. The area affected is governed by the aquifer's transmissivity
(its permeability integrated over the saturated thickness) together with
its storativity, the time elapsed and the position of any recharge boundary,
rather than by the tunnel's dimensions. The resulting settlement trough may be
an order of magnitude wider than the volume-loss trough of
Chapter~\ref{ch:settlement-damage} (Figure~\ref{fig:drawdown-trough}), is slower to develop, and is frequently
attributed to the tunnel long after construction has finished. Where
compressible layers overlie a dewatered aquifer, this mechanism can dominate
the total settlement.

\begin{figure}[htbp]
\centering
\datafigure[\textwidth]{ch07/drawdown-trough}
\caption{Normalised settlement troughs (schematic parameters): volume loss over a 20~m deep tunnel ($i = 0.5z_0$) against settlement taken proportional to drawdown from a drained tunnel, $\mathrm{erfc}\bigl(x/2\sqrt{Dt}\bigr)$ with $D = T/S = 2\times10^{-3}$~m\textsuperscript{2}/s. The drawdown trough is six times wider within a month and keeps widening for years.}
\label{fig:drawdown-trough}
\end{figure}

The environmental consequences usually govern. Groundwater-dependent
ecosystems, licensed bores, baseflow to watercourses and saline intrusion are
all sensitive to drawdown, and limits on each are commonly written into
approval conditions. Those limits then act as design constraints: a drawdown
condition determines the lining drainage philosophy as directly as any
structural calculation.

Control measures span a wide range of permeability and head
(Table~\ref{tab:groundwater-control}). Dewatering suits
permeable ground where drawdown is acceptable; cut-off walls and
pre-excavation grouting reduce inflow where it is not; ground freezing handles
high heads and fine soils at high cost. Compressed air working, once the usual
means of holding back water in soft ground, is now used sparingly because of
the decompression risk to the workforce, and pressurised closed-face machines
have largely taken over its role (Section~\ref{sec:guidance-interventions}).
Selection is treated in
Chapter~\ref{ch:ground-improvement}. Control measures are chosen against the
hydrogeological model, and a misdiagnosed flow regime produces a measure that
does not work: grouting a rock mass whose flow is carried by two discrete
features, for instance.

\begin{table}[htbp]
\caption{Groundwater control measures and the conditions each suits. Selection is treated in Chapter~\ref{ch:ground-improvement}.}
\label{tab:groundwater-control}
\centering
\small
\begin{tabularx}{\textwidth}{>{\raggedright\arraybackslash}p{0.26\textwidth} Y Y}
\toprule
\textbf{Measure} & \textbf{Where it is used} & \textbf{Limitation or remark} \\
\midrule
Dewatering & Permeable ground where drawdown is acceptable & Lowers the water table, with the settlement and environmental effects of this section \\
Cut-off walls & Where drawdown is not acceptable & Reduce inflow \\
Pre-excavation grouting & Where drawdown is not acceptable & Reduces inflow; does not work where the flow regime has been misdiagnosed \\
Ground freezing & High heads and fine soils & High cost \\
Compressed air working & Soft ground; now used sparingly & Decompression risk to the workforce \\
Pressurised closed-face machines & Soft ground under water pressure & Have largely taken over the role of compressed air \\
\bottomrule
\end{tabularx}
\end{table}

\subsection{Inrush and water chemistry}
\label{sec:inrush-chemistry}

Sudden inflow is a different hazard from steady seepage. It arises where the
excavation intersects a confined high-head feature: a karst conduit, a
fault-hosted aquifer, a buried valley, or old flooded workings. The volumes
can overwhelm pumping capacity within minutes, and the consequence is loss of
the heading and, on occasion, of life.

The standard defence is procedural rather than analytical. Probe drilling
ahead of the face, maintained as a policy through reaches identified as
at-risk in the geological model, provides warning and allows controlled
bleed-off or pre-excavation grouting before the feature is opened. Probing
costs standstill time at the face, which can be estimated in advance, whereas
the loss from opening an unprobed feature cannot be bounded in advance.

Water chemistry determines both what the structure must survive and what may
be discharged. Sulfate, chloride, low pH and
dissolved carbon dioxide all attack concrete, and the exposure classification
that sets binder type and cover comes from groundwater analysis carried out
during investigation (Chapter~\ref{ch:linings}). Sulfide-bearing rock exposed
to air generates acid rock drainage that can persist for decades in spoil
stockpiles. Discharge is regulated on quality as well as volume, and treatment
requirements discovered late are expensive; the chemistry should be
established while the investigation is still in the field.

\section{Summary}
\label{sec:hydrogeology-summary}

This chapter separated the effects of groundwater on strength, on inflow and
on lining load, and left face stability to Section~\ref{sec:face-support}.
The flow regime, porous or fracture-controlled, has to be diagnosed before
inflow is calculated, and an inflow prediction is stated as a range because
measured permeability depends on the scale of the test. The choice between a
drained and a tanked lining follows from the head, the permeability and the
acceptability of drawdown, and it is made early. Chapter~\ref{ch:in-situ-stress}
turns to the stress field in which the excavation is made. The control
measures are selected in Chapter~\ref{ch:ground-improvement}, the lining is
designed in Chapter~\ref{ch:linings}, and settlement is assessed in
Chapter~\ref{ch:settlement-damage}.

\chapterreferences
